\section{The 2024 sBootstrap and the SO(32) completion}
\label{sec:so32completion}
The subsequently supplied paper by Rivero~\cite{Rivero2024} makes an
important refinement to the preceding comparison. Sections 5 and 8
explicitly distinguish its proposed scalar composites from ordinary
QCD-scale mesons and diquarks, and envisage a different high-energy
binding mechanism. Thus the distinctive hypothesis is that a subset
of observed SM fermions supplies the constituents of scalar partners;
it is not necessarily an identification of those partners with known
QCD hadrons. This narrows the contrast with new strong dynamics in
Refs.~\cite{ARS2025,AGRS2026}, although the constituent and operator
dictionaries remain different.

\paragraph{What the five-flavor count supplies.}
Let $V=(5,3)$ under $SU(5)_f\times SU(3)_c$.
Its complex dimension is 15, so its underlying real space has dimension
30. The natural orthogonal extension is therefore $SO(30)$.
Its adjoint already contains the channels used in the scalar count:
\begin{align}
435&=224+105+\overline{105}+1,\\
224&=(24,1)+(1,8)+(24,8),\\
105&=(15,\bar3)+(10,6).
\end{align}
Indeed Sections 3.2 and 7.1 of Ref.~\cite{Rivero2024} explicitly
recognize $SO(30)$ as sufficient for the relevant decompositions.

The minimal extension to 32 real labels is one additional complex
color--flavor singlet:
\begin{equation}
16=(5,3)+(1,1),\qquad
32_{\mathbb C}=(5,3)+(\bar5,\bar3)+1+1.
\end{equation}
These are internal representation dimensions, not counts of spinor
components. A sterile-neutrino-like endpoint and its conjugate are a
possible physical interpretation, not a result of the count. Here the
fifteen labels are five quark flavors times three colors, rather than
the fifteen Weyl fermions of one SM family. The familiar right-handed
neutrino completion of a different family embedding therefore does
not automatically apply. Adding the omitted colored top instead
gives 18 complex labels and points to $SO(36)$.

\paragraph{The extra adjoint content.}
Taking the exterior square of $\mathbb R^{30}\oplus\mathbb R^2$
gives the branching under $SO(30)\times SO(2)$,
\begin{equation}
496=(435,1)+(30,2)+(1,1).
\end{equation}
The difference is 61 adjoint components, not just two extra scalar
states. In color--flavor language the added content is
$2(5,3)+2(\bar5,\bar3)+(1,1)$, precisely the mixed channels and
additional singlet visible in Eq.~(1) of Ref.~\cite{Rivero2024}.
A candidate neutral endpoint must account for these mixed colored
channels. Its electric charge also needs an explicit embedding.
All these identities are dimensionless and make no convention
setting physical $\hbar$ to one.

\paragraph{An enlargement is not yet a selection rule.}
The paper suggests replacing color $SU(3)$ by $U(3)$.
However, the unitary action $A\otimes B$ of
$SU(5)\times U(3)$ on $\mathbb C^{15}$ already acts orthogonally
on $\mathbb R^{30}$, with real determinant
$|\det_{\mathbb C}(A\otimes B)|^2=1$.
Its image, up to a finite kernel, lies in $U(15)\subset SO(30)$.
The extra central phase alone consequently does not force 32
real directions. Nor does adding singlets by itself demonstrate
interactions invariant under the full larger algebra.

One possible additional selector is the standard ten-dimensional
$N=1$ supergravity/super-Yang--Mills Green--Schwarz anomaly-cancellation
condition. Its gravitational anomaly requires 496 gauge generators;
within the orthogonal family,
\begin{equation}
\frac{n(n-1)}{2}=496\quad\Longrightarrow\quad n=32.
\end{equation}
Gauge and mixed anomaly factorization are also required; the type-I
string implements a related selection through tadpole cancellation.
This is a conditional higher-dimensional input, not a consequence of
the four-dimensional scalar count. The presence of 496 scalar
components alone does not supply the necessary ten-dimensional
gauge sector, and $E_8\times E_8$ is another familiar string possibility.

The missing ingredient is therefore a physical reason for exactly
one complex neutral sector and for the full symmetry, or a genuine
string-consistency condition selecting it. A neutral constituent is
a useful guess; no particular particle assignment, electroweak
completion or binding dynamics is derived here.
